% Options for packages loaded elsewhere
\PassOptionsToPackage{unicode}{hyperref}
\PassOptionsToPackage{hyphens}{url}
%
\documentclass[
]{article}
\usepackage{amsmath,amssymb}
\usepackage{lmodern}
\usepackage{iftex}
\ifPDFTeX
  \usepackage[T1]{fontenc}
  \usepackage[utf8]{inputenc}
  \usepackage{textcomp} % provide euro and other symbols
\else % if luatex or xetex
  \usepackage{unicode-math}
  \defaultfontfeatures{Scale=MatchLowercase}
  \defaultfontfeatures[\rmfamily]{Ligatures=TeX,Scale=1}
\fi
% Use upquote if available, for straight quotes in verbatim environments
\IfFileExists{upquote.sty}{\usepackage{upquote}}{}
\IfFileExists{microtype.sty}{% use microtype if available
  \usepackage[]{microtype}
  \UseMicrotypeSet[protrusion]{basicmath} % disable protrusion for tt fonts
}{}
\makeatletter
\@ifundefined{KOMAClassName}{% if non-KOMA class
  \IfFileExists{parskip.sty}{%
    \usepackage{parskip}
  }{% else
    \setlength{\parindent}{0pt}
    \setlength{\parskip}{6pt plus 2pt minus 1pt}}
}{% if KOMA class
  \KOMAoptions{parskip=half}}
\makeatother
\usepackage{xcolor}
\usepackage{longtable,booktabs,array}
\usepackage{calc} % for calculating minipage widths
% Correct order of tables after \paragraph or \subparagraph
\usepackage{etoolbox}
\makeatletter
\patchcmd\longtable{\par}{\if@noskipsec\mbox{}\fi\par}{}{}
\makeatother
% Allow footnotes in longtable head/foot
\IfFileExists{footnotehyper.sty}{\usepackage{footnotehyper}}{\usepackage{footnote}}
\makesavenoteenv{longtable}
\setlength{\emergencystretch}{3em} % prevent overfull lines
\providecommand{\tightlist}{%
  \setlength{\itemsep}{0pt}\setlength{\parskip}{0pt}}
\setcounter{secnumdepth}{-\maxdimen} % remove section numbering
\ifLuaTeX
  \usepackage{selnolig}  % disable illegal ligatures
\fi
\IfFileExists{bookmark.sty}{\usepackage{bookmark}}{\usepackage{hyperref}}
\IfFileExists{xurl.sty}{\usepackage{xurl}}{} % add URL line breaks if available
\urlstyle{same} % disable monospaced font for URLs
\hypersetup{
  hidelinks,
  pdfcreator={LaTeX via pandoc}}

\author{}
\date{}

\begin{document}

\textbf{S4D-TAM: Semantic 4D Token Attention Map for Autonomous
Navigation of Unmanned Aerial Vehicles in GNSS-Degraded Environments}

\textbf{Dr. Mateusz Pomianek}

\emph{Department of Computer and Control Engineering,}

\emph{Rzeszów University of Technology, Poland, m.pomianek@prz.edu.pl}

\textbf{Abstract}

We propose S4D-TAM (Semantic 4D Token Attention Map), a novel navigation
system architecture for unmanned aerial vehicles (UAVs) operating in
GNSS-degraded environments. Unlike classical methods that repeatedly
analyze raw image streams, the proposed system transforms sensor data
(RGB, thermal, IMU, GNSS) into a persistent map of semantic 4D tokens.
Each token represents a spatial fragment or object instance and contains
its geometry, semantic embedding, dynamics, observation history, and
calibrated uncertainty. A hierarchical spatiotemporal transformer
updates the world model, aligns it with reference maps, and forecasts
occupancy in partially observable regions. The predictive map is then
used by a~planner that accounts for risk, energy cost, and information
value of future observations. The proposed representation aims to reduce
image processing redundancy, increase robustness to GNSS degradation,
and enable decision-making based on persistent spatial memory rather
than solely on the current sensor field of view. The article presents
the system concept, literature review, mathematical description, an
executable auditable reference implementation, and a two-level
experimental methodology separating external system comparison from
internal mechanism ablation. Confirmatory experimental results are not
yet reported.

\textbf{Keywords}: unmanned aerial vehicles, semantic SLAM, 4D tokens,
spatiotemporal transformer, autonomous navigation, GNSS degradation,
world model, occupancy forecasting

\textbf{Table of Contents}

1. Introduction

2. Literature Review

\begin{quote}
2.1. Semantic SLAM for UAVs

2.2. Point Cloud Transformers

2.3. 4D Occupancy Forecasting

2.4. World Models and Planning

2.5. UAV Navigation in GNSS-Denied Environments
\end{quote}

3. Research Gap and Motivation

4. S4D-TAM System Concept

\begin{quote}
4.1. Definition of the 4D Token

4.2. System Architecture

4.3. Current Reference Implementation and Target Research Modules
\end{quote}

5. Mathematical Description

\begin{quote}
5.1. Sensor Layer and Data Fusion

5.2. SLAM Front-End

5.3. Point Cloud Tokenization

5.4. Spatiotemporal Attention Mechanism

5.5. Object Memory and Semantics

5.6. Reference Map Alignment

5.7. Occupancy and Risk Forecasting

5.8. Trajectory Planning

5.9. Active Perception

5.10. Metrological Error Analysis of the S4D-TAM System

5.11. Extended Semantic Token Model

5.12. Bayesian Inference about Future Scene States

5.13. Bayesian Risk Estimation and Trajectory Planning

5.14. Adaptive World Model with Hierarchical Resolution

5.15. Computational Resource Management and Bottleneck Elimination
\end{quote}

6. Research Hypotheses

7. Proposed Experimental Validation

\begin{quote}
7.1. Datasets

7.2. Two-Level Comparison Design

7.2.1. External System Comparison

7.2.2. Internal Mechanism Study

7.3. Metrics

7.4. Confirmatory Protocol and Statistical Analysis
\end{quote}

8. Discussion

9. Conclusions and Future Research Directions

\textbf{1. Introduction}

Autonomous navigation of unmanned aerial vehicles (UAVs) in environments
with limited GNSS signal availability is one of the central challenges
of modern mobile robotics. Satellite signal degradation can occur in
dense urban environments, under heavy forest canopy, in tunnels and
canyons, and under deliberate jamming or spoofing conditions {[}19{]},
{[}20{]}, {[}47{]}, {[}48{]}. In such scenarios, a UAV must rely on
alternative information sources: RGB and thermal cameras, inertial
measurement units (IMU), altimeters, LiDARs, and reference maps
{[}26{]}, {[}53{]}, {[}54{]}, {[}57{]}.

Contemporary autonomous navigation systems perform pose estimation, map
construction, semantic segmentation, object recognition, occupancy
forecasting, and trajectory planning as separate, loosely coupled
modules {[}9{]}, {[}31{]}. This approach leads to computational
redundancy - the same images are repeatedly analyzed by different
modules, and semantic and geometric information is not efficiently
shared between perception and planning stages {[}12{]}, {[}40{]}.

In this work, we propose S4D-TAM (Semantic 4D Token Attention Map), a
novel navigation system architecture in which the fundamental
environment representation is a persistent, semantic, continuously
updated map of 4D tokens derived from the SLAM point cloud. RGB and
thermal camera data are used primarily to update the geometry and
semantic meaning of points and objects stored in the map. Planning, risk
assessment, and decision-making modules do not repeatedly analyze the
same images - they operate on a compressed spatial representation
containing geometry, semantics, dynamics, observation history, and
uncertainty {[}1{]}.

The key scientific novelty of the proposed system is the introduction of
a persistent memory map composed of semantic 4D tokens, which serves as
a unified representation for localization, environment understanding,
future scene state prediction, risk estimation, and UAV trajectory
planning. The system remembers not only where an object was located, but
also when and how many times it was observed, from which sensors the
information originated, how reliable its position and recognition are,
whether the object is static, dynamic, or periodically displaced, what
spatial relations hold between it and other terrain elements, what
invisible areas may be hidden behind it, and what significance the
object has for flight safety and planning {[}1{]}.

The article combines a conceptual system architecture with an executable
research reference and a reproducible benchmark framework. The current
repository implements token state and lifecycle management, proposal and
association, multimodal encoder interfaces, attention-related
processing, uncertainty calibration utilities, reference-map and
topological support, causal occupancy and motion forecasting,
deterministic risk-aware planning, telemetry, and the integrated
S4DTAMReference execution path. The final learned hierarchical model,
full public-dataset reproduction, and confirmatory results remain under
development. Section 7 therefore specifies a prospective two-level
validation protocol rather than reporting final performance results.

\textbf{2. Literature Review}

\textbf{2.1. Semantic SLAM for UAVs}

Semantic SLAM for aerial platforms must be lightweight, robust to GPS
signal loss, and capable of producing useful semantic maps for planning
and multi-robot coordination. Liu et al. proposed SlideSLAM - a sparse,
lightweight, decentralized metric-semantic SLAM for multi-robot
navigation that uses an object-based representation and enables
efficient map merging in heterogeneous aerial-ground teams {[}2{]}. Tao
et al. integrated active exploration with metric-semantic SLAM for
resource-constrained aerial robots, explicitly balancing exploration and
localization uncertainty (including semantic loop closures), which
informs S4D-TAM strategies for uncertainty-aware mapping in GNSS-denied
environments {[}3{]}. Liu et al. presented a real-time semantic SLAM
that detects and models tree trunks and planes from LiDAR to constrain
poses and enable drift compensation during long-range autonomous flights
under forest canopy, directly relevant to drift mitigation and semantic
constraints in S4D-TAM trajectory planning {[}4{]}.

Tian et al. proposed SAME - a ground-air collaboration system that fuses
semantic Octree units from a UAV and a UGV into a coherent global map
supporting collaborative exploration, offering a semantic octree
architecture and fusion approach useful in S4D-TAM for cross-platform
mapping and multi-perspective scene reconstruction {[}6{]}. Xiang and Su
developed a LiDAR-based semantic mapping pipeline that creates
hierarchical semantic maps together with occupancy grids for navigation
and task adaptation, illustrating practical semantic mapping designs and
hierarchical semantics that S4D-TAM can exploit for task-oriented scene
understanding {[}7{]}.

Li et al. introduced Hi-SLAM - a hierarchical categorical representation
within 3D Gaussian Splatting for compact encoding of multiple semantic
classes and improved semantic mapping and tracking, directly applicable
to S4D-TAM requirements for scalable semantic scene representations
{[}5{]}, {[}46{]}. More recent semantic SLAM works employing 3D Gaussian
Splatting include SGS-SLAM {[}27{]}, GS4 {[}28{]}, SemGauss-SLAM
{[}29{]}, SAP-SLAM {[}30{]}, DynaGS-SLAM {[}32{]}, GSFF-SLAM {[}33{]},
OVO-SLAM {[}34{]}, MSGS-SLAM {[}35{]}, NIS-SLAM {[}36{]}, and NEDS-SLAM
{[}45{]}, demonstrating growing interest in real-time dense semantic
representations for SLAM.

\textbf{2.2. Point Cloud Transformers}

Point cloud transformer architectures and related models provide
expressive spatiotemporal encoding of 3D geometry and dynamics, useful
for the point-cloud-based perception and temporal prediction modules of
S4D-TAM. PCPNet introduces a transformer network for future LiDAR point
cloud prediction with a semantic auxiliary training objective for
semantically consistent forecasts, directly applicable to S4D-TAM for
point-level prediction and semantic alignment {[}16{]}. Silva et al.
proposed S2TPVFormer, which extends TPV embeddings with hybrid
cross-view temporal attention for temporally consistent semantic 3D
occupancy, offering a spatiotemporal transformer design that S4D-TAM can
adopt for temporally consistent semantic occupancy mapping {[}12{]}.

Chen et al. presented AdaOcc, which employs adaptive forward-view
transformations and flow modeling (with strong Swin camera backbone
models) to improve voxel occupancy and flow prediction, providing
practical camera-to-3D transformer strategies applicable in
vision-oriented S4D-TAM modules {[}15{]}. Work on point cloud
forecasting as a proxy for 4D occupancy prediction redefines the problem
toward sensor-agnostic spatiotemporal (4D) occupancy forecasting,
directing S4D-TAM toward world-oriented 4D occupancy targets for robust
temporal modeling {[}8{]}.

Further relevant works include Point-BERT {[}43{]}, which introduces
masked point modeling for pre-training point cloud transformers; SAT3D
{[}42{]} for semantic segmentation of 3D scenes using slot attention;
and YOGO {[}37{]} and Efficient Point Transformer {[}38{]}, {[}39{]},
demonstrating efficient point cloud processing with token representation
and relation inference modules relevant to S4D-TAM token compression.

\textbf{2.3. 4D Occupancy Forecasting}

Accurate 4D occupancy forecasting provides the spatiotemporal scene
model required by S4D-TAM for risk assessment and trajectory planning.
Yu et al. proposed Drive-OccWorld, which adapts a vision-centric 4D
world model with memory normalization and action-conditioned decoding
for occupancy and flow prediction, and integrates forecasting with
end-to-end planning, providing an approach that S4D-TAM can leverage to
couple occupancy prediction with trajectory selection {[}8{]}. Zheng et
al. presented OccWorld - a spatiotemporal generative transformer
architecture for 4D occupancy prediction, demonstrating effectiveness on
driving datasets and offering S4D-TAM a transformer backbone for
long-horizon volumetric prediction {[}3{]}.

Ma et al. provided Cam4DOcc - a standardized camera-only benchmark and
four baselines for current and future occupancy estimation, providing
S4D-TAM with evaluation protocols and camera-only baseline designs for
4D occupancy forecasting {[}4{]}. Guo et al. proposed FSF-Net, which
uses coarse BEV scene flow plus vector-quantized latent features and
U-Net fusion to improve occupancy prediction accuracy, showing a
practical pipeline that S4D-TAM can adapt for efficient flow-conditioned
occupancy prediction {[}7{]}.

Wang et al. introduced OccSora, which uses a 4D scene tokenizer and a
diffusion transformer for generative long-horizon 4D occupancy
simulation conditioned on trajectories, offering S4D-TAM a generative
world simulator approach for scenario generation and planning evaluation
{[}10{]}. Liao et al. proposed I2-World with hierarchical intra/inter
tokenizers and an encoder-decoder design for efficient compression and
generation of temporally consistent 4D occupancy, presenting a
tokenization strategy that S4D-TAM can apply for compact world
representations and fast inference {[}14{]}.

Zheng et al. presented GaussianAD, which represents scenes through
semantic 3D Gaussians refined from images, predicts 3D flows for those
Gaussians, and integrates forecasting with planning - this sparse
semantic Gaussian approach is directly relevant to
S4D-TAM\textquotesingle s lightweight scene abstractions and
prediction-based planning {[}17{]}.

\textbf{2.4. World Models and Planning}

World models unify perception, prediction, and planning; for S4D-TAM,
they provide a~mechanism for simulating the future, evaluating
trajectories, and incorporating action-conditioned forecasts. Chen et
al. proposed DriveWorld - a Memory State-Space Model with dynamic memory
and static scene propagation for spatiotemporal pre-training,
demonstrating downstream gains in detection, mapping, tracking,
forecasting, and planning - relevant to S4D-TAM pre-training strategies
and multi-task world modeling {[}14{]}, {[}18{]}.

Zhao et al. used a contrastive world model objective to learn visual
representations for drone navigation and visuomotor policy learning,
illustrating how world model pre-training can improve navigation
policies, which S4D-TAM can leverage for improved visual generalization
{[}13{]}. Mei et al. proposed IR-WM, which focuses on predicting
residual changes relative to a~temporal BEV prior and couples
forecasting with planning, offering S4D-TAM a compact residual
prediction approach that reduces redundant static background modeling
and improves planning integration {[}10{]}.

Cai et al. presented NEUSIS - a compositional neuro-symbolic framework
combining neuro-symbolic perception, a probabilistic world model with
Bayesian filtering, and hierarchical search planning, exemplifying an
interpretable and uncertainty-aware world modeling and planning workflow
useful for S4D-TAM mission-level decision-making {[}11{]}. Liu et al.
proposed HPP, which integrates occupancy-conditioned prediction with
game-theoretic motion forecasting and planning, demonstrating methods
for aligning prediction patterns with planning objectives that S4D-TAM
can adopt for improved trajectory selection under uncertainty and
interaction {[}11{]}.

Feng et al. provided a comprehensive survey of world models for
autonomous driving, covering 4D occupancy forecasting, generative
simulation, and planning integration, providing a~conceptual foundation
and taxonomy that can guide S4D-TAM design choices in world model
selection, pre-training, and evaluation {[}9{]}.

\textbf{2.5. UAV Navigation in GNSS-Denied Environments}

Autonomous UAV navigation in GNSS-degraded environments requires
multi-sensor fusion, robust SLAM algorithms, and uncertainty-aware
planning strategies. Yue et al. proposed an air-ground cooperative UAV
autonomous navigation system for GPS-denied environments that combines
visual odometry, IMU, and communication with a ground robot for robust
localization {[}19{]}. Radwan et al. presented UAV-assisted visual SLAM
generating reconstructed 3D scene graphs in GPS-denied environments,
demonstrating scene graph representations for scene understanding and
planning {[}20{]}.

Nieuwenhuisen et al. developed an autonomous navigation system for micro
aerial vehicles in complex GNSS-denied environments, using LiDAR, IMU,
and obstacle-aware planning {[}26{]}, {[}53{]}. Yue et al. proposed
semantically guided visual autonomous navigation for UAVs that employs
semantic segmentation to identify navigable areas and obstacles for
improved path planning {[}54{]}. Hu et al. provided a review of
vision-based UAV localization and optimization-based path planning in
GPS-denied environments, synthesizing SLAM, object detection, and
trajectory planning methods {[}57{]}.

Further relevant works include UAV navigation systems using multi-sensor
fusion (IMU, camera, LiDAR, radar) {[}48{]}, {[}50{]}, {[}51{]},
{[}52{]}, {[}58{]}, KAZE feature-based SLAM {[}49{]}, map-aided
navigation and dead reckoning {[}56{]}, and platforms for
search-and-rescue missions {[}55{]}. Castano presented autonomous UAV
navigation in GPS-denied environments using LiDAR point clouds for
localization and mapping {[}59{]}.

\textbf{3. Research Gap and Motivation}

Contemporary solutions perform pose estimation, map construction,
semantic segmentation, object recognition, occupancy prediction,
trajectory planning, and risk assessment as separate, independent
modules {[}9{]}, {[}31{]}, {[}40{]}. Systems such as LVI-SAM integrate
LiDAR, camera, and IMU in a factor graph, while Kimera constructs
metric-semantic 3D models. Neither, however, constitutes a complete,
learned world model that exploits persistent map tokens for predicting
unseen hazards and guiding sensor attention {[}12{]}.

OpenScene enables assigning language-image embedding space features to
point cloud points. Open3DSG extends such approaches with inter-object
relationships, while Octree-Graph combines semantics, occupancy, and
graph structure. These solutions, however, primarily address scene
understanding rather than closed-loop UAV control operating under
limited computational resources and partial observability {[}6{]},
{[}7{]}.

Research on 4D occupancy forecasting has demonstrated that predicting
the future state of space can serve as an effective planning foundation
{[}8{]}, {[}10{]}, {[}14{]}, {[}17{]}. Nevertheless, automotive
applications, dense voxel representations, and full sensor sequence
processing still dominate. The proposed S4D-TAM system fills the gap
between semantic SLAM, 4D world models, and risk-aware planning for UAVs
in GNSS-degraded environments.

The key motivation is to transition from a classical geometric map to a
predictive spatial memory of the UAV - one that not only records where
objects are, but also predicts where they may appear, which areas are
unobservable, and what risk they pose to flight safety. This approach
can be termed a SLAM-native world model: a world model whose primary
representation space is the SLAM map rather than the image stream
{[}1{]}.

\textbf{4. S4D-TAM System Concept}

\textbf{4.1. Definition of the 4D Token}

A 3D token would represent a spatial fragment or object at a given
position:

\emph{tᵢ = {[}xᵢ, yᵢ, zᵢ{]}}

It would not, however, capture environmental changes or observation
history. In an autonomous system it is essential to distinguish a
building, tree, or rock from a person, vehicle, or object that may
change its position {[}1{]}.

A 4D token should therefore incorporate the time dimension and
additional attributes:

\emph{Tᵢ(t) = {[}pᵢ, Σᵢ, gᵢ, sᵢ, vᵢ, hᵢ, uᵢ, rᵢ{]}}

where the individual elements denote:

\begin{quote}
1. pᵢ ∈ ℝ³ - position in space

2. Σᵢ ∈ ℝ³ × 3 - position covariance matrix (geometric uncertainty)

3. gᵢ ∈ ℝdg - geometric descriptor (normals, curvature, local features)

4. sᵢ ∈ ℝds - semantic embedding (features from a neural network or
open-vocabulary model)

5. vᵢ ∈ ℝ³ - estimated velocity/motion vector

6. hᵢ ∈ ℝdh - observation history (observation count, timestamps, sensor
sources)

7. uᵢ ∈ {[}0,1{]} - overall token uncertainty (scalar reliability
measure)

8. rᵢ ∈ {[}0,1{]} - local risk indicator (hazard probability)
\end{quote}

The fourth dimension does not merely add time as an additional
coordinate. The token also stores persistence, observation ordering,
rate of change, and the predicted future state of the object {[}1{]},
{[}17{]}, {[}22{]}, {[}23{]}, {[}25{]}.

\textbf{4.2. System Architecture}

The S4D-TAM system consists of the following layers (Figure 1):

Layer 1: Sensor Fusion Integration of data from RGB and thermal cameras,
IMU, GNSS, altimeter, optionally LiDAR/radar, and reference maps. Each
measurement receives a~timestamp, coordinate frame transformation, and a
reliability measure {[}1{]}.

Layer 2: SLAM Front-End A lightweight visual front-end for feature
tracking, depth estimation, characteristic object detection, and feature
vector extraction. The geometric SLAM creates a~local point cloud and
estimates the UAV trajectory. The factor graph includes visual odometry,
inertial, GNSS, altimeter, semantic loop-closure, and reference-map
alignment factors {[}2{]}, {[}3{]}, {[}4{]}.

Layer 3: Tokenization and Hierarchical Representation The full point
cloud is divided into superpoints, object segments, adaptive octree
cells, local surface patches, and dynamic instances. Each segment is
encoded as one or more tokens. The hierarchical representation stores
large, regular surfaces at low resolution while edges, obstacles, and
dynamic objects are retained at higher resolution {[}5{]}, {[}6{]},
{[}7{]}, {[}37{]}, {[}38{]}.

Layer 4: Spatiotemporal Transformer A hierarchical attention mechanism
with local scope (tokens within a given spatial radius, predicted flight
corridor, sensor visibility region) and global scope (landmarks, major
terrain obstacles, dynamic objects, high-uncertainty areas,
elevated-risk locations). The attention weight depends on token features
and their mutual positions {[}8{]}, {[}12{]}, {[}14{]}, {[}16{]}.

Layer 5: Object Memory and Semantics Upon detecting a characteristic
object, the system creates an instance token storing a class probability
distribution, observation history, open-vocabulary embedding, and
spatial relations with other objects {[}5{]}, {[}11{]}, {[}34{]}.

Layer 6: Reference Map Alignment The system maintains two
representations: the current map Mt built on the fly, and the reference
map Mref derived from elevation maps, previous flights, or a simplified
terrain point cloud. Alignment uses a three-step procedure: global
topological and terrain descriptor matching, landmark token matching,
and local geometric optimization {[}4{]}, {[}19{]}, {[}20{]}.

Layer 7: Occupancy and Risk Forecasting The model predicts future space
occupancy and the probability of risk in partially observable areas. The
system infers that beyond a building edge lies an unobservable region, a
road or terrain trail may lead to an unseen object, a previously
detected dynamic object may have moved into a neighboring zone, and
terrain topography limits available escape trajectories {[}8{]},
{[}10{]}, {[}14{]}, {[}17{]}, {[}21{]}.

Layer 8: Trajectory Planning The planner minimizes an objective function
that accounts for trajectory length, energy cost, collision probability,
predicted environmental risk, and the cost arising from map uncertainty.
The system selects among: continuing flight, altitude change, heading
change, temporary hover in a safe zone, return along the memorized
route, limiting non-critical measurements, and increasing observation
intensity in high-uncertainty areas {[}11{]}, {[}19{]}, {[}41{]}.

Layer 9: Active Perception The active perception module selects camera
orientation, processing frequency, and sensory modality based on
expected uncertainty reduction relative to cost. The IR camera may be
activated more frequently in low illumination or in areas where prior
data indicates the possible presence of a thermal object. The RGB camera
need not be analyzed by the full network at every frame {[}1{]},
{[}3{]}, {[}6{]}.

\textbf{4.3. Current Reference Implementation and Target Research
Modules}

The S4D-TAM concept is accompanied by an executable Python research
reference. The implementation is intentionally modular so that
individual mechanisms can be tested, replaced, and ablated without
changing the evaluator boundary.

Implemented reference components include persistent token state and
lifecycle management (token.py, memory.py), token proposal and data
association (proposal.py, association.py), modality-oriented encoder
interfaces and masked fusion (encoders/), attention-related processing
(attention.py), calibration and uncertainty utilities (calibration.py),
reference-map and topological support (reference\_map.py, topology.py),
causal probabilistic occupancy and motion forecasting (forecasting.py),
deterministic risk-, energy-, time-, progress-, and information-aware
planning (planner.py), structured event telemetry (telemetry.py), and
the integrated S4DTAMReference pipeline (pipeline.py).

The reference implementation favors causality, deterministic replay,
explicit uncertainty, and auditability over maximum throughput. This is
advantageous during early numerical validation because an observed
difference can be traced to a defined mechanism rather than hidden
implementation nondeterminism.

Several formulations in Sections 5.12-5.15 describe target research
extensions rather than already completed backends. These include a
learned hierarchical spatiotemporal model, a DBN/particle or equivalent
multimodal probabilistic scene-state backend, an MPPI planner backend,
adaptive octree/LOD storage, a closed-loop active-perception controller,
and a safety-aware real-time scheduler. They are maintained as explicit
implementation milestones and must preserve the same causal and
evaluator-facing contracts when introduced.

\textbf{5. Mathematical Description}

\textbf{5.1. Sensor Layer and Data Fusion}

The system uses a sensor set 𝒮= \textbackslash s₁, s₂, \ldots,
sn\textbackslash, where each sensor sᵢ provides a measurement zᵢ(t) at
time t with uncertainty Σᵢ(t). Data fusion is performed in a factor
graph:

\emph{Ftotal = FVO + FIMU + FGNSS + Falt + Fsem + Fref}

where:

\begin{quote}
• FVO - visual odometry factor

• FIMU - inertial factor

• FGNSS - GNSS factor (with adaptive weight under degradation)

• Falt - altimeter factor

• Fsem - semantic loop-closure factor

• Fref - reference map alignment factor
\end{quote}

Each factor takes the form:

\emph{Fᵢ = ∑ₖ ‖ eᵢk ‖\_Σᵢk²}

where eᵢk is the residual error and Σᵢk is the covariance matrix
{[}2{]}, {[}3{]}, {[}4{]}.

\textbf{5.2. SLAM Front-End}

The SLAM front-end estimates the UAV trajectory X = \textbackslash x₁,
x₂, \ldots, xT\textbackslash{} and the point map M = \textbackslash m₁,
m₂, \ldots, mN\textbackslash{} by minimizing the cost function:

\emph{X*, M* = arg minX, M ∑ᵢ,j ρ( ‖ zᵢⱼ - h(xᵢ, mⱼ) ‖\_Σ\_ij² )}

where h(xᵢ, mⱼ) is the projection function of point mⱼ onto pose xᵢ, and
ρ(·) is a robust function (e.g., Huber) {[}2{]}, {[}4{]}.

The reconstruction scale is corrected using:

\emph{s* = arg mins ∑ₖ wₖ ‖ s · dₖVO - dₖref ‖²}

where dₖVO is the distance from visual odometry, dₖref is the reference
distance (from IMU, altimeter, GNSS, known object sizes, elevation map),
and wₖ are reliability weights {[}1{]}, {[}4{]}.

\textbf{5.3. Point Cloud Tokenization}

The point cloud M is divided into K tokens \textbackslash T₁, T₂,
\ldots, TK\textbackslash{} through hierarchical segmentation:

\emph{M = ⋃ᵢ=1K 𝒫ᵢ}

where 𝒫ᵢ is the set of points assigned to token Tᵢ. Segmentation may
use:

\begin{quote}
1. Superpoints - grouping of points with similar geometry and semantics
{[}37{]}, {[}38{]}

2. Object segments - instance detection {[}5{]}, {[}34{]}

3. Adaptive octree - hierarchical space partitioning {[}6{]}, {[}7{]}

4. Local surface patches - planes, edges, corners {[}4{]}

5. Dynamic instances - moving objects {[}17{]}, {[}22{]}
\end{quote}

For each token Tᵢ the following attributes are computed:

\emph{pᵢ = (1)/(\textbar 𝒫ᵢ\textbar) ∑m ∈ 𝒫ᵢ m}

\emph{Σᵢ = (1)/(\textbar 𝒫ᵢ\textbar) ∑m ∈ 𝒫ᵢ (m - pᵢ)(m - pᵢ)ᵀ}

\emph{gᵢ = GeomEncoder(𝒫ᵢ)}

\emph{sᵢ = SemEncoder(ℐᵢ, 𝒫ᵢ)}

where ℐᵢ is the set of images observing points 𝒫ᵢ {[}5{]}, {[}16{]},
{[}38{]}.

\textbf{5.4. Spatiotemporal Attention Mechanism}

The hierarchical spatiotemporal transformer updates token
representations through the attention mechanism:

\emph{Tᵢ\textquotesingle{} = Tᵢ + Attention(Qᵢ, K, V)}

where:

\emph{Qᵢ = WQ Tᵢ, quad K = WK T, quad V = WV T}

Local attention covers tokens within radius rlocal:

\emph{𝒩ᵢlocal = : ‖pᵢ - pⱼ‖ \textless{} rlocal\textbackslash{}}

Global attention covers tokens of particular significance:

\emph{𝒩ᵢglobal = : landmark(j) ∨ dynamic(j) ∨ uⱼ \textgreater{} θu ∨ rⱼ
\textgreater{} θr\textbackslash{}}

The attention weight accounts for spatial, temporal, and uncertainty
relations:

\emph{αᵢⱼ = (exp( fracQᵢᵀ Kⱼ)/(√dₖ) · φ(pᵢ, pⱼ, tᵢ, tⱼ) · ψ(uᵢ, uⱼ) )∑ₖ
∈ 𝒩ᵢ exp( (Qᵢᵀ Kₖ)/(√dₖ) · φ(pᵢ, pₖ, tᵢ, tₖ) · ψ(uᵢ, uₖ) )}

where φ(·) encodes spatial and temporal relations, and ψ(·) limits the
influence of uncertain observations {[}8{]}, {[}12{]}, {[}14{]},
{[}16{]}.

\textbf{5.5. Object Memory and Semantics}

Upon detecting a characteristic object, the system creates an instance
token:

\emph{Oₖ = {[}pₖ, Σₖ, bₖ, sₖ, P(cₖ \textbar{} z₁:t), hₖ{]}}

where bₖ is the bounding box, and P(cₖ \textbar{} z₁:t) is the
probability distribution over the object class cₖ given the observation
history z₁:t.

The class distribution is updated via Bayesian filtering:

\emph{P(cₖ \textbar{} z₁:t) = fracP(zt \textbar{} cₖ) P(cₖ \textbar{}
z₁:t-1)∑c\textquotesingle{} P(zt \textbar{} c\textquotesingle)
P(c\textquotesingle{} \textbar{} z₁:t-1)}

The open-vocabulary approach allows searching the map for objects and
properties that were not constrained to a closed class set during
training {[}5{]}, {[}11{]}, {[}34{]}.

\textbf{5.6. Reference Map Alignment}

The system maintains two representations: the current map Mt and the
reference map Mref. Alignment uses a three-step procedure:

Step 1: Global Alignment Determining the transformation Tglobal by
matching global descriptors:

\emph{Tglobal = arg minT ∑ᵢ ‖ dᵢt - d\_σ(i)ref ‖²}

where dᵢt and d\_σ(i)ref are global descriptors (e.g., topological,
terrain-based) {[}4{]}, {[}19{]}, {[}20{]}.

Step 2: Landmark Matching Refining the transformation by matching
landmark tokens:

\emph{Tlandmark = arg minT ∑ⱼ ∈ ℒ ‖ T pⱼt - p\_σ(j)ref ‖\_Σⱼ²}

where ℒ is the landmark set {[}2{]}, {[}4{]}.

Step 3: Local Geometric Optimization Final refinement via ICP or a
semantic variant:

\emph{Tfinal = arg minT ∑ᵢ,j wᵢⱼ ‖ T pᵢt - pⱼref ‖²}

where wᵢⱼ accounts for semantic consistency {[}4{]}, {[}5{]}.

The difference between the current map and the reference map indicates:
a localization error, an environmental change, appearance of a new
object, absence of a previously existing object, or unreliability of the
current GNSS measurement {[}1{]}, {[}19{]}.

\textbf{5.7. Occupancy and Risk Forecasting}

The system predicts future space occupancy and risk probability in
partially observable regions:

\emph{O(t+Δ t) = fpred(T₁:K(t), a(t))}

where O(t+Δ t) is the predicted occupancy at time t+Δ t, T₁:K(t) are
tokens at time t, a(t) is the planned action (trajectory), and fpred is
the predictive model (e.g., temporal transformer, diffusion model)
{[}8{]}, {[}10{]}, {[}14{]}, {[}17{]}.

The risk map is determined as:

\emph{R(p, t+Δ t) = P(Hp, t+Δ t \textbar{} M₁:t, Mref, Ct)}

where Hp, t+Δ t denotes a hazardous event at position p at time t+Δ t,
and Ct is the motion and environment context {[}1{]}, {[}21{]}.

The system infers risk in unobservable areas via:

\emph{Roccluded(p) = ∑c ∈ 𝒞 P(c \textbar{} Tvisible) · P(H \textbar{} c,
geometry(p))}

where 𝒞 is the set of contexts (road, terrain trail, building edge),
Tvisible are the visible tokens, and geometry(p) describes the terrain
topography at position p {[}1{]}, {[}11{]}, {[}21{]}.

\textbf{5.8. Trajectory Planning}

The planner minimizes the objective function:

\emph{J(ξ) = ∫₀ᵀ {[} L(ξ(t)) + E(ξ(t), ξ(t)) + λH PH(ξ(t)) + λR R(ξ(t),
t) + λU U(ξ(t)) {]} dt}

where:

\begin{quote}
• ξ(t) is the UAV trajectory

• L(ξ(t)) is the trajectory length

• E(ξ(t), ξ(t)) is the energy cost

• PH(ξ(t)) is the collision probability

• R(ξ(t), t) is the predicted environmental risk

• U(ξ(t)) is the cost arising from map uncertainty

• λH, λR, λU are weighting coefficients
\end{quote}

Collision probability is determined as:

\emph{PH(ξ(t)) = 1 - ∏ᵢ (1 - P(collision \textbar{} ξ(t), Tᵢ(t)))}

The uncertainty cost accounts for missing information:

\emph{U(ξ(t)) = ∑ᵢ ∈ 𝒩(ξ(t)) uᵢ · exp(-(‖ξ(t) - pᵢ‖²)/(2σ²))}

The system selects the action:

\emph{a* = arg mina ∈ 𝒜 J(ξ(a))}

where 𝒜 is the set of possible actions: continuing flight, altitude
change, heading change, temporary hover in a safe zone, return along the
memorized route, limiting non-critical measurements, and increasing
observation intensity in high-uncertainty areas {[}1{]}, {[}11{]},
{[}19{]},~{[}41{]}.

\textbf{5.9. Active Perception}

The active perception module selects a sensory action atsense that
maximizes expected uncertainty reduction relative to cost:

\emph{atsense = arg maxa ∈ 𝒜\_sense (I(T; z \textbar{} a))/(C(a))}

where I(T; z \textbar{} a) is the mutual information between tokens T
and measurement z given action a, and C(a) is the energy and
computational cost of the action {[}1{]}, {[}3{]}, {[}6{]}.

Mutual information is approximated as:

\emph{I(T; z \textbar{} a) ≈ H(T) - Ez\textbar a{[}H(T\textbar z){]}}

where H(T) is the token entropy before measurement and H(T\textbar z) is
the entropy after measurement.

The IR camera may be activated more frequently in low illumination or in
regions where prior data indicate the possible presence of a thermal
object. The RGB camera need not be processed by the full network at
every frame {[}1{]}.

\textbf{5.10. Metrological Error Analysis of the S4D-TAM System}

Metrological analysis forms the foundation for designing reliable
autonomous navigation systems. The objective of this subsection is to
formalize the error models of individual S4D-TAM system layers, derive
lower bounds on estimation uncertainty, and determine the cumulative
positioning error budget of the UAV in a GNSS-degraded environment.

\textbf{5.10.1. Sensor Layer Error Model}

Each sensor sᵢ ∈ 𝒮 is characterized by its own error model. Adopting the
general measurement form:

\emph{zᵢ(t) = hᵢ(x(t)) + nᵢ(t) + bᵢ(t)}

where hᵢ(·) is the observation function, nᵢ(t) sim 𝒩(0, Rᵢ) is the
measurement noise with covariance matrix Rᵢ, and bᵢ(t) is the systematic
error (bias).

IMU (gyroscope + accelerometer): The error model of the gyroscope and
accelerometer in an IMU extended with bias is described as:

\emph{ω(t) = ω(t) + bg(t) + ng(t), quad ng sim 𝒩(0, σg² I₃)}

\emph{a(t) = a(t) + ba(t) + na(t), quad na sim 𝒩(0, σa² I₃)}

where biases evolve according to a random walk model:

\emph{bg = nbg, quad nbg sim 𝒩(0, σbg² I₃)}

\emph{ba = nba, quad nba sim 𝒩(0, σba² I₃)}

The position error resulting from double integration of IMU grows over
time as:

\emph{σIMU(t) = (1)/(2)σa t² + (1)/(6)σba t³}

GNSS under degradation conditions: In GNSS-degraded environments we
introduce a signal quality indicator qGNSS(t) ∈ {[}0, 1{]}, where qGNSS
= 1 denotes full constellation coverage and qGNSS = 0 denotes complete
signal loss. The GNSS measurement covariance scales adaptively:

\emph{RGNSS(t) = fracRGNSS\^{}(0)qGNSS(t) + ε}

where RGNSS\^{}(0) is the baseline covariance at full signal and ε ll 1
is a regularization parameter preventing singularity. As qGNSS to 0, the
factor FGNSS is effectively removed from the factor graph through the
increase in its variance.

Visual odometry (VO) - relative drift: The visual odometry error
accumulates quasi-randomly, proportional to the traveled distance d:

\emph{σVO(d) = αVO d + βVO √(d)}

where αVO is the linear error coefficient (typically
0.5--2\textbackslash\% per meter traveled), and βVO describes the
diffusive component associated with stochastic visual noise.

Barometric altimeter: The altitude measurement error is dominated by
slow-varying temperature and ambient pressure drift:

\emph{zbaro(t) = ztrue(t) + bbaro + nbaro(t)}

where bbaro sim 𝒰(-δb, +δb) is the bias with a uniform distribution over
range δb ≈ 1--3 m, and nbaro sim 𝒩(0, σbaro²) with σbaro ≈ 0.1--0.5 m.

\textbf{5.10.2. Error Propagation in the Factor Graph}

Let ξ = {[}ptop, vtop, φtop, bgtop, batop{]}top ∈ ℝ¹⁵ be the navigation
state vector (p - position, v - velocity, φ - Euler angles, bg, ba - IMU
biases). The lower bound on the variance of an unbiased estimator ξ is
given by the Cramér-Rao Lower Bound (CRLB):

\emph{Cov(ξ) ≥ J-¹(ξ)}

where J(ξ) is the Fisher Information Matrix (FIM), the sum of
contributions from all factor graph factors:

\emph{J(ξ) = ∑ᵢ Hᵢtop Rᵢ-¹ Hᵢ + Jprior}

where Hᵢ = (∂ hᵢ)/(∂ ξ) is the Jacobian matrix of the observation
function of the i-th factor, and Jprior represents the prior information
contribution. The marginal position variance σp² is the diagonal
element:

\emph{σp² ≥ {[}J-¹(ξ){]}pp}

Navigation state covariance propagation follows the Riccati equation:

\emph{P(t) = F(t) P(t) + P(t) F(t)top + G(t) Q(t) G(t)top}

where F(t) is the inertial system state transition matrix, G(t) is the
noise input matrix, and Q(t) = diag(σa²I₃, σg²I₃, σba²I₃, σbg²I₃) is the
noise spectral density matrix.

When merging measurements from M sensors in the factor graph, the
combined FIM grows additively, reducing the estimator covariance:

\emph{Jfused = JIMU + JVO + qGNSS JGNSS + Jalt + Jsem + Jref}

where each term is weighted according to the reliability of the
corresponding source.

\textbf{5.10.3. Quantization and Discretization Error of Tokens}

The 3D space is divided into tokens of side length ell meters (§5.3).
The quantization operation mapping the coordinates of point p ∈ ℝ³ to
the centroid of token cₖ introduces a quantization error:

\emph{εq = p - cₖ, quad ‖εq‖ ≤ fracell√(3)2}

For a regular cubic grid of side length ell, the quantization error
variance along each axis is:

\emph{σq,1D² = (ell²)/(12)}

and the total 3D variance:

\emph{σq² = (ell²)/(4)}

An additional semantic error arises from imperfections of the
open-vocabulary classifier. Defining the classifier confusion matrix C ∈
ℝC × C (C classes), the semantic error of token k~is expressed by the
conditional entropy of the class assignment:

\emph{εsem\^{}(k) = -∑c=1C p(c\textbar Tₖ) log p(c\textbar Tₖ) =
H(c\textbar Tₖ)}

A value of εsem\^{}(k) = 0 indicates certain assignment, while
εsem\^{}(k) = log C indicates maximum semantic uncertainty. A token with
high εsem\^{}(k) receives increased covariance in the Fsem factor graph
factor.

\textbf{5.10.4. Lower Bound on Localization Error}

The total UAV positioning error is a combination of: visual odometry
drift σVO, uncertainty arising from IMU integration σIMU, GNSS error
σGNSS, and the contribution from semantic anchoring σsem. Under joint
optimization in the factor graph, the CRLB on the position error takes
the form:

\emph{σp,min²(t) = {[}Jfused-¹(t){]}pp = (1)/(∑ᵢ {[}Hᵢtop Rᵢ\^{}-1)
Hᵢ{]}pp}

In the scenario of complete GNSS absence (qGNSS = 0), position error
grows according to a~modified drift model:

\emph{σp²(t) = σp,0² + αVO² d(t)² + γsem Nloop-¹}

where σp,0² is the initial variance (determined by GNSS before
degradation), d(t) is the total distance traveled, Nloop is the number
of successful semantic loop closures, and γsem is the semantic
correction coefficient. Each semantic loop closure reduces σp²
proportionally to the strength of the semantic evidence, determined by
the Fisher information of factor Fsem.

\textbf{5.10.5. Occupancy Forecasting Error}

The 4D occupancy predictor (§5.7) produces a probabilistic map O(x, t+Δ
t) based on the token state at time t. Prediction error grows with the
temporal horizon Δ t. Modeling the evolution of token Tₖ as a
linear-Gaussian stochastic process with transition matrix Φ and noise
diffusion Qdyn:

\emph{Tₖ(t+Δ t) = Φ(Δ t) Tₖ(t) + ωₖ, quad ωₖ sim 𝒩(0, Qdyn(Δ t))}

the prediction error covariance is:

\emph{Σpred(Δ t) = Φ(Δ t) Σₖ(t) Φ(Δ t)top + Qdyn(Δ t)}

For static objects Φ ≈ I and Qdyn ≈ 0, while for dynamic objects Qdyn(Δ
t) ∝ Δ t. This results in a linear growth of prediction uncertainty with
horizon:

\emph{σpred²(Δ t) = σpred,0² + σdyn² · Δ t}

The binary occupancy classification error is defined by the Brier Score:

\emph{BS(Δ t) = (1)/(\textbar 𝒱\textbar) ∑x ∈ 𝒱 (O(x, t+Δ t) - Otrue(x,
t+Δ t))²}

where 𝒱 is the set of workspace voxels. The expected value of BS(Δ t) is
lower bounded by the entropy of the dynamic environment.

\textbf{5.10.6. Cumulative Error Budget}

The table below summarizes the error sources of the S4D-TAM system with
their typical values for the target configuration (lightweight UAV in an
urban environment with partial GNSS degradation):

\begin{longtable}[]{@{}
  >{\raggedright\arraybackslash}p{(\columnwidth - 6\tabcolsep) * \real{0.2657}}
  >{\raggedright\arraybackslash}p{(\columnwidth - 6\tabcolsep) * \real{0.2343}}
  >{\raggedright\arraybackslash}p{(\columnwidth - 6\tabcolsep) * \real{0.2188}}
  >{\raggedright\arraybackslash}p{(\columnwidth - 6\tabcolsep) * \real{0.2812}}@{}}
\toprule()
\begin{minipage}[b]{\linewidth}\raggedright
\textbf{Error Source}
\end{minipage} & \begin{minipage}[b]{\linewidth}\raggedright
\textbf{Characteristic}
\end{minipage} & \begin{minipage}[b]{\linewidth}\raggedright
\textbf{Typical value (1σ)}
\end{minipage} & \begin{minipage}[b]{\linewidth}\raggedright
\textbf{Time/Distance dependence}
\end{minipage} \\
\midrule()
\endhead
IMU noise (accelerometer) & White noise σa & 0.05 m/s² & σ ∝ t²
(integration) \\
IMU bias drift & Random walk σba & 0.002 m/s²/√s & σ ∝ t³/2 \\
Visual odometry & Proportional drift αVO & 1\textbackslash\% per meter &
σ ∝ d \\
GNSS (full signal) & σGNSS\^{}(0) & 1.5 m (horizontal) & Constant \\
GNSS (degraded) & σGNSS(q) & 5--50 m & ∝ 1/qGNSS \\
Barometric altimeter & σbaro + bias & 0.3 m + 2 m & Thermal drift \\
Token quantization & σq & ell/2 (e.g., 0.5 m for ell = 1 m) & Constant
(depends on ell) \\
Semantic error & Classifier entropy \$H(c &
\textbackslash mathbf\{T\})\$ & 0.1--0.8 bit \\
Occupancy prediction (Δ t = 1 s) & BS(Δ t) & 0.05--0.15 & ∝ Δ t \\
Total position error (IMU+VO only, d = 100 m) & CRLB σp & ≈ 2--5 m & ∝
d \\
Total position error (IMU+VO+sem, Nloop = 5) & CRLB σp & ≈ 0.5--1.5 m &
Bounded by loop closures \\
\bottomrule()
\end{longtable}

The error budget indicates that the key mechanism for limiting position
drift under GNSS-degraded conditions is the effectiveness of semantic
loop closures. Each loop closure with high semantic confidence (low
εsem\^{}(k)) reduces position variance proportionally to the Fisher
information of factor Fsem, enabling maintenance of the error below the
operational threshold σp \textless{} 2 m even when flying through zones
of complete GNSS loss extending to several kilometers.

\textbf{5.11. Extended Semantic Token Model}

The basic geometric token introduced in §5.3 is replaced by an extended
semantic-temporal token 𝒯ₖ, constituting a complete profile of an object
recognized in the UAV\textquotesingle s environment. Each token stores
not only the object\textquotesingle s position, but a complete record of
the multidimensional observation history:

\emph{𝒯ₖ = ⟨ μₖ, Σₖ, pc\^{}(k), 𝒪ₖ, δₖ, vₖ, ℛₖ, Ωₖ, ρₖ ⟩}

where the individual components have the following meaning:

\begin{quote}
• μₖ ∈ ℝ³ - estimated position (centroid) in the global frame

• Σₖ ∈ ℝ³ × 3 - position covariance matrix (localization uncertainty)

• pc\^{}(k) = (cⱼ \textbar{} 𝒯ₖ)\textbackslash ⱼ=1C - semantic class
probability distribution

• 𝒪ₖ = \textbackslash(tᵢ, sᵢ, wᵢ, zᵢ)\textbackslash ᵢ=1nₖ - observation
history: time, sensor ID, reliability weight, measurement

• δₖ ∈ , dynamic, periodic\textbackslash{} - motion state classification

• vₖ ∈ ℝ³ - estimated velocity vector (for δₖ ≠ static)

• ℛₖ = \textbackslash rₖl\textbackslash{} - set of spatial-semantic
relations to neighboring tokens

• Ωₖ - occlusion region (areas hidden by object k)

• ρₖ ∈ {[}0,1{]} - significance coefficient for flight safety and
planning
\end{quote}

Weighted multi-sensor fusion:

Each observation in the history 𝒪ₖ carries a reliability weight
depending on the sensor source, detection confidence, and time elapsed
since the last observation:

\emph{wᵢ = wbase(sᵢ)· confᵢ · fdecay(t - tᵢ)}

where wbase(sᵢ) is the baseline reliability of sensor sᵢ (calibrated a
priori), confᵢ ∈ {[}0,1{]} is the detection confidence, and fdecay is
the temporal decay function:

\emph{fdecay(Δ t) = exp(-λs · Δ t)}

The constant λs is a sensor-specific parameter: RGB cameras decay faster
(λRGB \textgreater{} λLiDAR) due to sensitivity to illumination changes.
The position estimate and its covariance are determined by weighted
information fusion:

\emph{Σₖ-¹ = ∑ᵢ=1nₖ wᵢ · Σᵢ-¹, qquad μₖ = Σₖ ∑ᵢ=1nₖ wᵢ · Σᵢ-¹ zᵢ\^{}(k)}

Bayesian update of the semantic class:

At each new observation znew, the token class is updated according to
Bayes\textquotesingle{} rule:

\emph{p(cⱼ \textbar{} 𝒯ₖ, znew) ∝ p(znew \textbar{} cⱼ) · p(cⱼ
\textbar{} 𝒯ₖ)}

The prior p(cⱼ \textbar{} 𝒯ₖ) comes from the previous iteration, and the
likelihood p(znew \textbar{} cⱼ) is the output of the open-vocabulary
classifier. The accumulation of observations causes semantic certainty
to increase monotonically with nₖ.

Motion state classification:

Based on the observation history 𝒪ₖ, the system classifies the token
into one of three dynamic classes. Define the sequence of estimated
displacements:

\emph{Δ μₖ\^{}(i) = μₖ\^{}(tᵢ) - μₖ\^{}(tᵢ₋₁)}

The classification criterion is based on a statistical test:

\emph{δₖ = begincases static \& if ‖Δ μₖ‖̄ \textless{} εs and σ\_Δ
\textless{} ε\_σ \textbackslash\textbackslash{} periodic \& if
FFT\textbackslash‖Δ μₖ\^{}(i)‖\textbackslash{} contains a dominant
frequency f* \textbackslash\textbackslash{} dynamic \& otherwise
endcases}

where εs and ε\_σ are experimentally calibrated thresholds.

Spatial relations between tokens:

A pair of tokens (𝒯ₖ, 𝒯l) is connected by a relation rₖl describing
their mutual geometric and semantic conditioning:

\emph{rₖl = ⟨ dₖl, eₖl, ψₖl, αₖl ⟩}

where dₖl = ‖μₖ - μl‖ is the distance, eₖl is the unit direction vector,
ψₖl is the ontological relation (e.g., above, adjacent, blocks), and αₖl
is the position uncertainty correlation coefficient. These relations
form a scene graph 𝒢 = (𝒱, ℰ), in which vertices are tokens and edges
carry relations rₖl.

Modeling of occluded areas:

For each token 𝒯ₖ and the current UAV position xUAV, the system
determines the occlusion cone Ωₖ(xUAV) as the set of points
geometrically hidden by the object:

\emph{Ωₖ(xUAV) = x ∈ ℝ³ : ∃ λ \textgreater{} 1, x = xUAV + λ (μₖ -
xUAV), ‖x - axisₖ‖ \textless{} rₖ}

where rₖ is the estimated object radius. These regions enrich the
uncertainty map with knowledge about the unknown surroundings.

Flight safety significance coefficient:

Token significance ρₖ is a function of the semantic class, observation
history, and operational context:

\emph{ρₖ = ρclass(cₖ*) · (1 - exp(-β\_ρ · nₖ)) · fcontext(𝒯ₖ, xUAV)}

where ρclass(cₖ*) is the baseline class significance (e.g., ρwire =
0.95, ρtree = 0.6, ρbuilding = 0.8), the factor (1 - exp(-β\_ρ nₖ))
reflects growing confidence with the number of observations, and
fcontext accounts for distance and current flight conditions.

\textbf{5.12. Bayesian Inference about Future Scene States}

Implementation status. The current executable reference uses a strictly
causal probabilistic occupancy and motion forecaster driven only by
observations available up to the prediction time. It produces Bernoulli
occupancy probabilities, flow means, uncertainty estimates,
observability masks, and physically timed forecast horizons, including
irregular or dropped samples. This deterministic causal baseline is
preferred for initial validation because future-frame leakage is
structurally impossible. The DBN and particle approximation below define
the target probabilistic extension for multimodal future-state
uncertainty and are not claimed as the current default backend.

Predicting future scene states is a key element of proactive UAV
navigation. The system models the scene as a Dynamic Bayesian Network
(DBN), in which the global scene state at time t is defined as:

\emph{St = 𝒯ₖ(t)ₖ=1K}

Scene state transition model:

Assuming conditional independence of tokens with respect to the common
environmental background, the transition model factorizes as:

\emph{p(St+Δ t \textbar{} St) = ∏ₖ=1K p(𝒯ₖ(t+Δ t) \textbar{} 𝒯ₖ(t), δₖ,
𝒢t)}

Separate transition models are applied for static, dynamic, and periodic
tokens:

\emph{p(μₖ(t+Δ t) \textbar{} μₖ(t), δₖ) = begincases 𝒩(μₖ, Σₖ + Qs Δ t)
\& δₖ = static \textbackslash4pt{]} 𝒩(μₖ + vₖ Δ t, Σₖ + Qd Δ t) \& δₖ =
dynamic \textbackslash4pt{]} ∑ⱼ αⱼ 𝒩(μₖ\^{}(j), Σₖ\^{}(j)) \& δₖ =
periodic endcases}

where Qs and Qd are process noise matrices for static and dynamic
objects, and the periodic model is a Gaussian mixture with components
μₖ\^{}(j) corresponding to typical positions within the cycle.

Particle approximation of the scene state distribution:

Due to the nonlinearity of models and multimodality of distributions,
the posterior probability of the scene state is approximated by a set of
Np particles \textbackslash(st\^{}(i), wt\^{}(i))\textbackslash ᵢ=1Np:

\emph{p(St \textbar{} z₁:t) ≈ ∑ᵢ=1Np wt\^{}(i) δ(St - st\^{}(i))}

The particle filter recursion proceeds in three steps:

Prediction:

\emph{st₊₁\^{}(i) sim p(St₊₁ \textbar{} st\^{}(i))}

Weight update:

\emph{wt₊₁\^{}(i) ∝ wt\^{}(i) · p(zt₊₁ \textbar{} st₊₁\^{}(i))}

Normalization and resampling:

\emph{wt₊₁\^{}(i) ← fracwt₊₁\^{}(i)∑ⱼ wt₊₁\^{}(j), qquad resample when
Neff = (1)/(∑ᵢ (wt\^{}(i)))² \textless{} Np / 2}

The predictive occupancy map at horizon Δ t is the particle-weighted
average:

\emph{O(x, t+Δ t) = ∑ᵢ=1Np wt\^{}(i) O(x \textbar{} st+Δ t\^{}(i))}

Inference about occluded areas:

For the occlusion region Ωₖ determined in §5.11, the occupancy
probability is inferred from the semantic context of the occluding token
and spatial environment structure statistics:

\emph{p(occupied(x) \textbar{} 𝒯ₖ) = p₀(cₖ*) + (1 - p₀(cₖ*)) exp(-(d(x,
∂Ωₖ)²)/(2σ\_Ω²)), quad x ∈ Ωₖ}

where p₀(cₖ) is the baseline probability of an obstacle existing behind
an object of class cₖ (e.g., p₀(building) = 0.8, p₀(tree) = 0.3), and
σ\_Ω controls the uncertainty range. This inference allows the planner
to account for hazards invisible to sensors.

Anomaly and event detection:

The deviation of the observed state Stobs from the predicted state St is
measured by the Mahalanobis distance:

\emph{DM(t) = √((Stobs) - St)top ΣS\^{}-1 (Stobs - St)}

Exceeding the threshold DM(t) \textgreater{} Dalarm triggers heightened
attention mode, in which the system increases processing frequency and
directs sensors toward the anomaly area.

\textbf{5.13. Bayesian Risk Estimation and Trajectory Planning}

Implementation status. The executable reference planner currently uses
deterministic finite-horizon beam search over admissible acceleration
controls with explicit kinematic limits and decomposed collision-risk,
energy, time, goal-progress, and information-value terms. This approach
is more auditable than stochastic sampling for early confirmatory
experiments because identical inputs generate the same candidate
evaluation and decision trace. MPPI remains a planned backend for richer
dynamics and will be introduced behind the same planner contract only
after direct comparison with the deterministic reference.

The S4D-TAM system integrates the risk map directly into the trajectory
planning process, treating collision risk as a probabilistic field
determined in real time from the current Bayesian state.

Risk field definition:

The risk field ℛ(x, t) at spatial point x at time t is the sum of
weighted contributions from all tokens, including occluded areas:

\emph{ℛ(x, t) = ∑ₖ=1K ρₖ · pocc(x, 𝒯ₖ, t) · pcoll(x, 𝒯ₖ)}

where the occupancy probability pocc comes from the Bayesian prediction
(§5.12), and the geometric collision probability is modeled as the
complementary CDF of the safety distance:

\emph{pcoll(x, 𝒯ₖ) = 1 - Φ(fracd(x, μₖ) - rsafe√(σₖ,r)² + σ\_UAV²)}

where rsafe is the minimum safety radius, σₖ,r is the object size
uncertainty, and σUAV accounts for UAV position uncertainty.

Risk in occluded areas is elevated proportionally to the uncertainty of
the occluded region:

\emph{ℛ(x, t) ← ℛ(x, t) + κ\_Ω · 1x ∈ ⋃ₖ Ωₖ · maxₖ p₀(cₖ*)}

Trajectory cost function:

The optimal trajectory τ* minimizes a multi-component cost function:

\emph{J(τ) = wL L(τ) + wE E(τ) + wR ∫\_τ ℛ(x(t), t) dt + wU ∫\_τ 𝒰(x(t),
t) dt - wI ℐ(τ)}

where:

\begin{quote}
• L(τ) - trajectory length

• E(τ) - energy cost (dependent on UAV dynamics)

• ℛ - integrated collision risk along the route

• 𝒰(x,t) = σp²(x,t) + κ\_Ω · 1x ∈ Ω - world model uncertainty

• ℐ(τ) - information gain (rewarding active exploration of
high-uncertainty zones)
\end{quote}

Information gain is expressed through the expected scene entropy
reduction:

\emph{ℐ(τ) = H(St) - E\_τ{[}H(St \textbar{} z₁:τ){]}}

Bayesian trajectory optimization (BTO):

\emph{τ* = arg min\_τ ∈ ℱ ES{[}J(τ \textbar{} S){]}}

For a future high-dimensional stochastic planning backend, direct
minimization of the expected value may be approximated by particle
trajectory sampling via MPPI (Model Predictive Path Integral). The
following formulation defines that target extension; it is not required
by the current deterministic reference planner:

\emph{τ* ≈ frac∑m=1M exp(-(1)/(λMPPI) J(τm)) τm∑m=1M exp(-(1)/(λMPPI)
J(τm))}

where M sampled trajectories \textbackslash τm\textbackslash{} are
generated randomly from a normal distribution around the nominal
trajectory, and λMPPI is a temperature parameter controlling the
exploration-exploitation tradeoff.

Multi-level decision system:

The risk field ℛ(x, t) defines four decision levels:

\emph{UAV Mode = begincases nominal \& ℛ(xUAV, t) \textless{} ℛsafe
\textbackslash\textbackslash{} alert \& ℛsafe ≤ ℛ \textless{} ℛalert
\textbackslash\textbackslash{} avoidance \& ℛalert ≤ ℛ \textless{}
ℛavoid \textbackslash\textbackslash{} emergency \& ℛ ≥ ℛavoid endcases}

In alert mode, the planner selects avoidance trajectories with wider
margins; in avoidance mode, an immediate maneuver is required; in
emergency mode, the UAV initiates a safe landing or return procedure.

\textbf{5.14. Adaptive World Model with Hierarchical Resolution}

Implementation status. The current token memory already enforces bounded
history and configurable limits on token count, map memory, and update
time, together with lifecycle activation, sleeping, reactivation,
merging, and removal rules. These mechanisms implement the
resource-bounded principle of this section. The full adaptive octree and
Level-of-Detail storage policy described below remains a planned
world-model backend and will be validated separately before becoming the
default representation.

Uniform tokenization resolution is a fundamental bottleneck in systems
operating under limited computational resources. The S4D-TAM system
introduces an Adaptive World Occupancy Model (AWOM) based on a
hierarchical octree structure.

Hierarchical octree structure:

The UAV\textquotesingle s operational space is represented by an
adaptive octree 𝒪 with node resolutions ell ∈ \textbackslash ell₁, ell₂,
\ldots, ellL\textbackslash, where ell₁ \textless{} ell₂ \textless{} ⋯
\textless{} ellL:

\emph{ell₁ = 0.1 m,quad ell₂ = 0.5 m,quad ell₃ = 2 m,quad ell₄ = 10 m}

Each octree node at level ellⱼ is compressed into a Level of Detail
(LOD) token representation at level j:

\emph{LODₖ ∈ \textbackslash0, 1, 2, 3\textbackslash{}}

\begin{quote}
• LOD 0: token stores only (μₖ, ρₖ, cₖ*) - centroid, significance, and
dominant class

• LOD 1: LOD 0 + approximate point cloud (up to 64 pts) + pc\^{}(k)

• LOD 2: LOD 1 + full attention embedding + observation history 𝒪ₖ

• LOD 3: LOD 2 + motion prediction model + spatial relations ℛₖ + region
Ωₖ
\end{quote}

Resolution adaptation criterion:

The optimal token resolution ellₖ is determined as a trade-off between
perception quality, computational resources, and location significance:

\emph{ellₖ = ell₁ · 2\^{} Lₖ, qquad Lₖ = leftlfloor Lmax · fdist(dₖ) ·
frisk(ρₖ)-¹ · fcomp(Cavail)-¹ rightrfloor}

where:

\emph{fdist(dₖ) = min(1, fracd(μₖ, xUAV)dmax), quad frisk(ρₖ) = ρₖ +
ε\_ρ, quad fcomp(Cavail) = fracCavailCnominal}

Interpretation: tokens close to the UAV, with high safety significance,
and under conditions of good resource availability receive higher
resolution; distant, low-significance tokens, or those during CPU
overload, receive lower resolution.

Dynamic processing priorities:

Each token 𝒯ₖ has a priority scalar πₖ computed at each planning cycle:

\emph{πₖ = α\_π · ρₖ + β\_π · (1)/(λmin)(Σₖ) + ε + γ\_π · fage(t -
tlast,k) + δ\_π · 1\_δₖ ≠ static}

where:

\begin{quote}
• ρₖ - safety significance (§5.11)

• λmin(Σₖ) - smallest eigenvalue of the covariance matrix (lower value =
high confidence, lower update priority)

• fage(t - tlast,k) = 1 - exp(-ν · (t - tlast,k)) - aging factor
encouraging refresh of stale tokens

• 1\_δₖ ≠ static - priority bonus for dynamic and periodic tokens
\end{quote}

Computational budget allocation across layers:

The total computational budget Ctotal is distributed among the
processing pipeline layers:

\emph{Ctotal = Csensor + CSLAM + Ctoken + Cattn + Cpred + Cplan}

Constraints: ∑ᵢ Cᵢ ≤ Ctotal and Csensor ≥ Csensormin (minimum resources
for perception safety). When Cavail \textless{} Cnominal, reduction
occurs in cascade from the layer with the lowest safety priority:

\emph{Cᵢ ← max(Cᵢmin, Cᵢ · fracCavailCnominal), qquad i ∈ , pred, plan,
token, SLAM, sensor\textbackslash{}}

\textbf{5.15. Computational Resource Management and Bottleneck
Elimination}

Implementation status. Structured telemetry and resource budgets already
expose update-time, token-count, memory, lifecycle, and decision events
needed for profiling and replay. Predictive queue-overflow control,
EDF-RT scheduling, guaranteed safety-critical CPU reservations, and the
full survival-mode controller described below remain planned real-time
systems extensions rather than features of the current Python reference.

Intelligent resource management is a necessary condition for realizing
all described S4D-TAM layers on embedded-class UAV hardware. This
section formalizes the data flow model, bottleneck identification, and
elimination mechanisms.

Processing pipeline throughput analysis:

The S4D-TAM pipeline consists of N sequential stages with processing
times Tᵢ and queue sizes qᵢ(t). System throughput is limited by the
bottleneck stage:

\emph{Θsys = (1)/(maxᵢ) Tᵢ, qquad i* = arg maxᵢ Tᵢ}

For the S4D-TAM pipeline with stages: \{acquisition, SLAM front-end,
tokenization, attention, Bayesian prediction, planner\}, the nominal
cycle time is:

\emph{Tcycle = ∑ᵢ=1⁶ Tᵢ + Tcomm}

where Tcomm accounts for data bus latencies between modules.

Bottleneck prediction:

The system monitors queue lengths qᵢ(t) and predicts overload via
Bayesian methods. Modeling task arrival as a Poisson process with
intensity λᵢ(t), the probability of queue overflow within horizon Δ t
is:

\emph{P(qᵢ(t+Δ t) \textgreater{} qmax) = 1 - ∑n=0qmax-qᵢ(t) ((λᵢ · Δ
t)n)/(n!) e\^{}-λᵢ Δ t}

When P \textgreater{} Palarm, the system proactively reduces the load of
stage i by:

\begin{quote}
• Reducing tokenization resolution (tokenization stage)

• Skipping updates of low-priority tokens πₖ (attention stage)

• Shortening the prediction horizon Δ tpred (prediction stage)

• Simplifying the trajectory sampler (reducing M in MPPI, planner stage)
\end{quote}

Adaptive processing frequency:

The effective processing frequency is dynamically scaled to available
resources:

\emph{fproc(t) = min(fnom, fracCavail(t)Cper)}

When fproc \textless{} fmin, the system automatically simplifies the
world model: it switches to survival mode with only a local attention
map within radius rsurvival and a reactive planner based on VFH+ (Vector
Field Histogram+).

Token memory management - TTL and LOD-compression policy:

Total token memory is bounded by the budget:

\emph{Mtokens = ∑ₖ=1K m(LODₖ) ≤ Mbudget}

where m(LODⱼ) is the token size at level j (e.g., m(0) = 64 B, m(1) =
512 B, m(2) = 4 KB, m(3) = 32 KB). The memory management policy combines
three mechanisms:

TTL (Time-To-Live): a token is deactivated if t - tlast,k \textgreater{}
TTTL and ρₖ \textless{} ρprune.

Merging: static tokens at distance dₖl \textless{} ellₖ with
KL-divergence DKL(pc\^{}(k) ‖ pc\^{}(l)) \textless{} εKL are merged into
a common token:

\emph{μₖl = (nₖ μₖ + nl μl)/(nₖ + nl), qquad Σₖl-¹ = Σₖ-¹ + Σl-¹}

LOD-compression: tokens beyond the active processing range are degraded
to a lower LOD, freeing operational memory.

Task scheduling with safety guarantee:

The task scheduler allocates CPU time to S4D-TAM system modules
according to an EDF-RT (Earliest Deadline First - Real-Time) policy
extended with a safety coefficient:

\emph{Priorityᵢ(t) = fracρᵢmaxTdeadline,i(t) - t + κsafety · 1ᵢ ∈
𝒞\_safety}

where ρᵢmax = maxₖ ρₖ in the task queue of stage i, Tdeadline,i(t) is
the absolute completion deadline, and 𝒞safety is the set of
safety-critical stages (sensor acquisition, SLAM front-end, collision
detection). Safety-critical tasks 𝒞safety have a guaranteed minimum CPU
time regardless of system load, ensuring minimally safe operation even
under full resource saturation.

\textbf{6. Research Hypotheses}

The confirmatory mechanism study uses seven preregistered hypotheses.
Each hypothesis compares the full S4D-TAM model with a variant that
disables exactly one mechanism while data, seeds, preprocessing,
optimization budget, and all other settings remain fixed. External
systems are deliberately excluded from this causal ablation family.

H1: Semantics improve navigation in scenarios containing dynamic
objects. The primary endpoint is mission success; collisions per
kilometer and semantic mIoU are secondary outcomes.

H2: Temporal state improves 1 s future-state prediction. The primary
endpoint is forecast occupancy IoU at 1 s; flow EPE and temporal flip
rate are secondary outcomes.

H3: Calibrated uncertainty reduces navigation risk. The primary endpoint
is collisions per kilometer; ECE, NLL, and near misses are secondary
outcomes.

H4: Topological support increases mission success. Path efficiency and
ATE RMSE are secondary outcomes.

H5: Reference-map support reduces localization error. The primary
endpoint is ATE RMSE; RPE and final drift percentage are secondary
outcomes.

H6: Predictive risk modeling reduces collisions. The primary endpoint is
collisions per kilometer; near misses and minimum clearance are
secondary outcomes.

H7: Token lifecycle management reduces computational cost without
materially degrading navigation quality. Confirmation requires both
lower p95 latency and non-inferior mission success with a preregistered
margin of 2 percentage points; map bytes and peak RSS are secondary
outcomes.

H1-H6 are superiority hypotheses in the preregistered direction. H7
combines efficiency superiority with mission-success non-inferiority.
The seven confirmatory decisions form one family and use the Holm
procedure defined in Section 7.4.

\textbf{7. Proposed Experimental Validation}

\textbf{7.1. Datasets}

The validation framework targets TartanAir, Blackbird, MARSIM, and
AeroVerse, each converted to the common SequenceData contract with
immutable source version, license record, coordinate convention,
calibration metadata, and SHA-256 manifest. The cohort identifier
frozen-2026-01 denotes the study cohort rather than a vendor release and
must resolve to exact source versions and file hashes.

Data are split at the environment or flight level, never by adjacent
frames. The train split is the only source of learned weight updates;
calibration is reserved for thresholds, normalization, early stopping,
and uncertainty calibration; test remains inaccessible until code,
artifacts, and analysis are frozen.

The preregistered offline mechanism study uses 20 preselected test
sequences per dataset, yielding 80 paired sequence-level units. Systems
validation additionally defines 30 paired SIL scenarios with five seeds,
20 HIL scenarios with five seeds, and 20 paired controlled real-flight
missions per variant when the safety gate remains open.

A dedicated UAV dataset remains desirable for synchronized RGB, thermal,
IMU, GNSS, point cloud, elevation/reference map, static and dynamic
labels, visibility or observability masks, and controlled
sensor-degradation scenarios. Metrics that require unavailable ground
truth are reported as unavailable rather than silently omitted or
imputed.

7.2. Two-Level Comparison Design

The evaluation is split into two non-interchangeable levels. Both levels
use common dataset and metric contracts where applicable, but they
answer different scientific questions and produce separate statistical
reports.

\textbf{7.2.1. External System Comparison}

The external comparison asks whether the complete S4D-TAM system is
competitive with independently implemented navigation and SLAM systems
under the same data, hardware policy, sensor availability, and evaluator
definitions.

The mandatory primary baseline set comprises ORB-SLAM3, VINS-Mono,
FAST-LIO2, and LIO-SAM. Each baseline is pinned to an upstream revision
and fixed environment or container, calibration, sensor-topic mapping,
loop-closure policy, warm-up policy, and hardware configuration. Results
enter the benchmark only through the common AlgorithmResult artifact
schema.

Additional external systems may be included only when their sensing
assumptions are compatible, their exact implementation can be
reproduced, and their output can be normalized without method-specific
post-processing. Such additions must be frozen before the external study
is executed.

This level measures complete-system competitiveness. A difference
relative to an external baseline cannot be interpreted as the causal
effect of an individual S4D-TAM component.

\textbf{7.2.2. Internal Mechanism Study}

The internal comparison tests mechanism contribution using the full
system and seven single-component ablations: H1\_no\_semantics,
H2\_no\_temporal\_state, H3\_no\_calibrated\_uncertainty,
H4\_no\_topology, H5\_no\_reference\_map, H6\_no\_risk\_prediction, and
H7\_no\_token\_lifecycle.

Each variant differs from full by exactly one switch and uses its own
frozen trained artifact. External systems never appear in this matrix.
This separation prevents confounding complete-system identity with the
causal effect of removing one S4D-TAM mechanism.

\textbf{7.3. Metrics}

Primary endpoints are deliberately compact and preregistered:
localization uses ATE RMSE after SE(3) alignment without scale;
forecasting uses occupancy IoU at 1 s and, for broader characterization,
3 s; safety uses mission success and collisions per kilometer;
efficiency uses p95 latency and persistent map-memory footprint.

Localization and SLAM secondary metrics include RPE translation RMSE,
final drift in meters and percentage of reference path length, and
rotation RPE when orientation ground truth is available. Relocalization
success rate, time to relocalize, and post-relocalization pose error are
reported when the dataset supports those events.

Semantic metrics include mIoU and macro F1. Forecasting metrics include
per-horizon occupancy IoU/F1, flow EPE, Brier score, NLL, and ECE. Pose
uncertainty is evaluated by NEES, 95\% coverage, and NLL where
covariance output is available. Risk prediction may additionally report
AUROC, false-alarm rate, and miss rate.

Navigation metrics include mission success, collisions per kilometer,
near misses, minimum clearance, and path efficiency. Efficiency metrics
include latency p95, map bytes, peak RSS, CPU time, and
implementation-specific telemetry such as processed-token counts when
consistently available.

Scale error, embedding quality, spatial-relation accuracy, emergency
replanning counts, FPS, and energy consumption may be reported as
secondary or exploratory metrics when their ground truth and measurement
procedure are explicitly defined. They do not become confirmatory
endpoints merely because an implementation emits them.

A missing metric caused by absent ground truth or absent algorithm
output is recorded as unavailable. Algorithm failure, timeout, abort, or
collision is not treated as a missing observation and remains visible in
the safety analysis.

\textbf{7.4. Confirmatory Protocol and Statistical Analysis}

The preregistered H1-H7 study uses sequence or mission level inference
rather than treating individual frames as independent samples.
Continuous outcomes are summarized with paired mean differences and 95\%
BCa bootstrap confidence intervals using 10,000 resamples. Mission
success is modeled as a paired scenario-level probability outcome, while
distance-normalized event counts use Poisson or negative-binomial
modeling with log distance as an offset according to the preregistered
overdispersion rule.

The seven H1-H7 confirmatory decisions form one family and are corrected
with the sequential Holm procedure at family-wise alpha = 0.05. H1-H6
require the corrected result and confidence interval to support the
predicted superiority direction. H7 additionally requires
mission-success non-inferiority within 2 percentage points while
improving p95 latency.

No missing technical result is imputed. Aborts and algorithm failures
remain in the denominator for safety outcomes; complete-case reporting
is accompanied by the preregistered worst-case sensitivity analysis.
Exclusions are limited to integrity failures, excessive missing
timestamps, independently confirmed sensor faults, or operational
constraints defined before result inspection.

The external system comparison is reported separately from this H1-H7
inferential family. It uses the same frozen datasets and common
evaluators, but its effect estimates describe differences between
complete systems rather than component-level causal effects.

\textbf{8. Discussion}

The proposed S4D-TAM system represents an attempt to transition from a
classical geometric map to a predictive spatial memory of the UAV. The
key advantage is amortized perception - the cost of image analysis is
incurred once, while the result remains available in spatial memory
across many subsequent planning cycles {[}1{]}. This approach may
substantially reduce computational redundancy in systems where the same
images are repeatedly analyzed by different modules {[}12{]}, {[}40{]}.

The hierarchical token representation allows storing large, regular
surfaces at low resolution, while edges, obstacles, and dynamic objects
are retained at higher resolution {[}5{]}, {[}6{]}, {[}7{]}, {[}37{]},
{[}38{]}. This is consistent with the direction set by Superpoint
Transformer, Octree-Graph, and other architectures demonstrating that
aggregating points into higher-level structures can substantially reduce
representation complexity.

The spatiotemporal attention mechanism enables modeling of long-range
relations between tokens while maintaining computational efficiency by
restricting attention to local neighborhoods and global landmarks
{[}8{]}, {[}12{]}, {[}14{]}, {[}16{]}. Token-uncertainty-weighted
attention limits the influence of uncertain observations on world model
updates, which may improve system stability under conditions of partial
observability {[}1{]}, {[}11{]}.

Occupancy and risk forecasting in partially observable areas constitutes
the most interesting element of the system - the transition from
detection to prediction {[}1{]}, {[}8{]}, {[}10{]}, {[}14{]}, {[}17{]},
{[}21{]}. The system does not claim to have recognized an object it has
not seen, but determines the probability of risk existing in partially
observable space based on geometric context, semantics, and observation
history. This approach may enable planning that is more defensive and
risk-aware than classical methods based solely on the current sensor
field of view.

Reference map alignment increases localization robustness to periodic
GNSS loss or~inconsistency {[}4{]}, {[}19{]}, {[}20{]}. The difference
between the current map and the reference map can indicate not only a
localization error, but also an environmental change, the appearance of
a new object, or the unreliability of the current GNSS measurement
{[}1{]}. The three-step alignment procedure (global descriptors,
landmarks, local geometric optimization) enables efficient map alignment
at varying density and resolution.

Active perception maximizing expected uncertainty reduction relative to
energy and computational cost may substantially improve system
efficiency under resource constraints {[}1{]}, {[}3{]}, {[}6{]}. The IR
camera may be activated more frequently in low illumination or in
regions where prior data indicate the possible presence of a thermal
object, while the RGB camera need not be processed by the full network
at every frame.

The limitations of the proposed approach include:

\begin{quote}
1. Implementation complexity - the system requires integration of
multiple components (SLAM, segmentation, transformer, prediction,
planning) and their efficient synchronization.

2. Computational requirements - despite amortized perception, the
spatiotemporal transformer and 4D occupancy forecasting may require
significant computational resources.

3. Reference map quality - the effectiveness of reference map alignment
depends on the availability and currency of elevation maps or previous
flights.

4. Uncertainty calibration - determining calibrated token uncertainty
and risk probability requires careful training and validation.

5. Generalization - the model may require fine-tuning for new
environments and operational conditions.
\end{quote}

Future work should focus on completing and training the learned S4D-TAM
model, validating public-dataset converters and pinned external
baselines, executing the preregistered H1-H7 mechanism study, and
progressing from offline evaluation through SIL, HIL, and controlled
real-flight validation. Additional engineering work includes the
probabilistic DBN/particle backend, MPPI planning, adaptive octree/LOD
storage, active perception, and safety-aware real-time scheduling.

\textbf{9. Conclusions and Future Research Directions}

This paper has proposed S4D-TAM (Semantic 4D Token Attention Map) - a
novel navigation system architecture for unmanned aerial vehicles
operating in GNSS-degraded environments. The key novelty is the
introduction of a persistent memory map composed of semantic 4D tokens,
which serves as a unified representation for localization, environment
understanding, future scene state prediction, risk estimation, and UAV
trajectory planning.

The main conceptual contributions include:

\begin{quote}
1. Definition of the 4D token combining geometry, semantics, time,
uncertainty, and observation history in a single compact representation.

2. Hierarchical system architecture integrating sensor fusion, SLAM,
tokenization, spatiotemporal transformer, object memory, reference map
alignment, occupancy and risk forecasting, trajectory planning, and
active perception.

3. Spatiotemporal attention mechanism accounting for spatial, temporal,
and token uncertainty relations.

4. Risk forecasting in partially observable areas based on geometric
context, semantics, and observation history.

5. Trajectory planning minimizing length, energy cost, collision
probability, predicted environmental risk, and the cost arising from map
uncertainty.

6. Active perception maximizing expected uncertainty reduction relative
to energy and computational cost.
\end{quote}

Seven preregistered mechanism hypotheses have been formulated for
semantics, temporal state, calibrated uncertainty, topology,
reference-map support, predictive risk, and token lifecycle. Their
causal full-vs-H1-H7 study is explicitly separated from the external
system comparison against ORB-SLAM3, VINS-Mono, FAST-LIO2, and LIO-SAM.

Future research directions include:

\begin{quote}
1. Completion and training of the learned S4D-TAM model followed by
two-level validation: external system comparison on TartanAir,
Blackbird, MARSIM, and AeroVerse, and the preregistered internal
full-vs-H1-H7 mechanism study.

2. Computational efficiency optimization through quantization, pruning,
knowledge distillation, and dedicated hardware accelerators.

3. Extension with additional sensor modalities (radar, ultrasound,
magnetometers) and operational scenarios (multi-robot coordination,
long-term missions, dynamic environments).

4. Integration with language models for natural-language-guided
navigation and high-level task understanding.

5. Ablation studies of individual system components to identify key
factors influencing performance.

6. Validation under real operational conditions on UAV platforms in
GNSS-degraded environments.
\end{quote}

The proposed S4D-TAM system represents a step toward autonomous UAVs
capable of navigating in complex, partially observable environments with
limited GNSS signal availability, by exploiting persistent predictive
spatial memory in place of repeated analysis of raw sensor streams.

\textbf{Acknowledgements}

The author thanks the reviewers for their valuable comments and the
Department of Computer Science and Automation of Rzeszów University of
Technology for research support.

\textbf{References}

\begin{quote}
\emph{Editorial note: Numbers {[}12{]}, {[}14{]}, {[}16{]}, {[}18{]}
contain numbering conflicts (the same number used for two different
works in different article sections) - these require manual confirmation
of full bibliographic data in the final version.}
\end{quote}

{[}2{]} Liu X., van der Lei J., Prabhu A., Tao Y., Spasojević I.,
"SlideSLAM: Sparse, Lightweight, Decentralized Metric-Semantic SLAM for
Multi-Robot Navigation," arXiv:2406.17249, 2024.

{[}3{]} Tao Y., Liu X., Spasojević I., Agarwal S., Kumar V., "3D Active
Metric-Semantic SLAM," IEEE Robot. Autom. Lett., vol. 9, 2024. DOI:
10.1109/lra.2024.3363542.

{[}4{]} Liu X., Nardari G.V., Ojeda F.C., Tao Y., Zhou A., Donnelly T.,
Qu C., Chen S.W., Taylor C.J., Kumar V., "Large-Scale Autonomous Flight
With Real-Time Semantic SLAM Under Dense Forest Canopy," IEEE Robot.
Autom. Lett., vol. 7, no. 2, 2022. DOI: 10.1109/lra.2022.3154047.

{[}5{]} Li B., Cai Z., Wu C., Lin P., Zhang Y., Rezatofighi H.,
"Hi-SLAM: Scaling-up Semantics in SLAM with a Hierarchically Categorical
Gaussian Splatting," arXiv:2409.12518, 2024.

{[}6{]} Tian X., Deng Y., Tang Y., Wang J., Dang R., "SAME: Ground-Air
Collaborative Semantic Active Mapping and Exploration," Proc. 2024 IEEE
Int. Conf. Unmanned Systems (ICUS), 2024. DOI:
10.1109/icus61736.2024.10839908.

{[}7{]} Xiang Z., Su J., "Towards Customizable Robotic Disinfection with
Structure-Aware Semantic Mapping," IEEE Access, vol. 9, 2021. DOI:
10.1109/ACCESS.2021.3062043.

{[}8{]} Yu Y., Mei J., Jiang X., Yang H., Xu Z., Luo Y., Wang Y., Yang
Z., Li Z., "Drive-OccWorld: Vision-Centric 4D Occupancy Forecasting and
Planning via World Models," Proc. AAAI Conf. Artificial Intelligence,
vol. 39, no. 9, 2025. DOI: 10.1609/aaai.v39i9.33010.

{[}9{]} Feng T., Wang W., Yang Y., "A Survey of World Models for
Autonomous Driving," arXiv:2501.11260, 2025.

{[}10{]} Mei J., Yu Y., Jiang X., Wang S., Li Z., "IR-WM: Vision-Centric
4D Occupancy Forecasting and Planning via Implicit Residual World
Models," arXiv:2510.16729, 2025.

{[}11{]} Cai Z., Rezatofighi H., Haffari G., Baktashmotlagh M., "NEUSIS:
A Compositional Neuro-Symbolic Framework for Autonomous Perception,
Reasoning, and Planning in Complex UAV Search Missions,"
arXiv:2409.10196, 2024.

{[}12{]} Rosen D.M., Carlone L., "Scene Representations for Robotic
Spatial Perception," Annual Review of Control, Robotics, and Autonomous
Systems, 2024. DOI: 10.1146/annurev-control-040423-030709. {[}Conflict:
same number used for S2TPVFormer (Silva et al.) in §2.2 - to be
verified.{]}

{[}13{]} Zhao J., Li C., Guo W., Zhao J., "Learning Visual
Representation for Autonomous Drone Navigation via a Contrastive World
Model," IEEE Trans. Artificial Intelligence, vol. 4, no. 6, 2023. DOI:
10.1109/tai.2023.3283488.

{[}14{]} Chen M., Zhao D., Xu S., Chen Y., Xu C., Liu J., "DriveWorld:
4D Pre-Trained Scene Understanding via World Models for Autonomous
Driving," arXiv:2405.04390, 2024. {[}Conflict: same number used for
I2-World in §2.3 - to be verified.{]}

{[}15{]} Chen D., Han W., Fang J., Shen J., "AdaOcc: Adaptive Forward
View Transformation and Flow Modeling for 3D Occupancy and Flow
Prediction," arXiv preprint, 2024.

{[}16{]} Liao Z., Zhong F., Zhang Z., Fu Y., Sun J., Li X., Cheng X.,
"I2-World: Intra-Inter Tokenization for Efficient Dynamic 4D Scene
Forecasting," arXiv:2507.09144, 2025. {[}Conflict: same number used for
PCPNet in §2.2 - to be verified.{]}

{[}17{]} Zheng W., Zhou B., Liu Y., Wang J., "GaussianAD:
Gaussian-Centric End-to-End Autonomous Driving," arXiv:2412.10371, 2024.

{[}18{]} Zheng W., Chen Y., Huang Z., Zhong Y., Lei W., Sun Z., Zhan J.,
Liu Y., "OccWorld: Learning a 3D Occupancy World Model for Autonomous
Driving," Proc. ECCV, Springer LNCS, 2024. {[}Conflict: same number used
for DriveWorld in §2.4 - to be verified.{]}

{[}19{]} Yue W., Nie Y., Song X., Zhao Z., Wang J., Shang Z.,
"Air-Ground Collaborative Navigation for UAVs in Complex Urban
Environments," Drones, vol. 9, no. 6, 2025. DOI: 10.3390/drones9060442.

{[}20{]} Radwan N., Olivares-Mendez M.A., Voos H., Burgard W.,
"UAV-Assisted Visual SLAM with Scene Graphs for Outdoor Navigation,"
Proc. 2024 Int. Conf. Unmanned Aircraft Systems (ICUAS), 2024. DOI:
10.1109/icuas60882.2024.10556948.

{[}21{]} Liu H., Zhao R., Ma Y., Liu X., Fu T., Zhao Z., Wang S., Chen
S., Ye M., Chen X., "Hybrid-Prediction Integrated Planning for
Autonomous Driving," IEEE Trans. Pattern Anal. Mach. Intell., 2025. DOI:
10.1109/tpami.2025.3526936.

{[}22{]} Wang L., Zheng W., Huang Z., Chen Y., Chen B., Li H., "OccSora:
4D Occupancy Generation Models as World Simulators for Autonomous
Driving," arXiv preprint, 2024.

{[}23{]} Ma J., Chen X., Deschaud J.E., Marcotegui B., Goulette F., Hu
X., "Cam4DOcc: Benchmark for Camera-Only 4D Occupancy Forecasting in
Autonomous Driving Scenes," arXiv preprint, 2023.

{[}24{]} Wang J., Gao P., Zhou X., Guo Z., "TartanAir: A Dataset to Push
the Limits of Visual SLAM," Proc. IEEE/RSJ Int. Conf. Intelligent Robots
and Systems (IROS), 2020.

{[}25{]} Chen J., Wei Z., Ruan K., Yang J., Zhang C., "Risk Occupancy: A
New and Efficient Paradigm through Vehicle-Road-Cloud Collaboration,"
arXiv:2408.07367, 2024.

{[}26{]} Nieuwenhuisen M., Droeschel D., Stückler J., Behnke S.,
"Autonomous Navigation for Micro Aerial Vehicles in Complex GNSS-Denied
Environments," Proc. IEEE Int. Symp. Safety, Security, and Rescue
Robotics (SSRR), 2015. DOI: 10.1109/SSRR.2015.7443012.

{[}27{]} Hu C., Liu N., Chen X., Lyu C., "SGS-SLAM: Semantic Gaussian
Splatting for Neural Dense SLAM," arXiv:2402.03246, 2024.

{[}28{]} Yugay V., Wang Y., Pollefeys M., Rezatofighi H., "GaussianSLAM:
Photo-realistic Dense SLAM with Gaussian Splatting," arXiv:2312.10070,
2023. {[}to be verified - cited in text as GS4; check association with
this work.{]}

{[}29{]} Zou Y., He Z., Xu B., Cheng H., "SemGauss-SLAM: Dense Semantic
Gaussian Splatting SLAM," arXiv:2403.07494, 2024.

{[}30{]} Huang H., Li L., Cheng H., Yeung S.K., "Photo-SLAM: Real-time
Simultaneous Localization And Photorealistic Mapping for Monocular,
Stereo, and RGB-D Cameras," Proc. IEEE/CVF CVPR, 2024. {[}to be verified
- cited in text as SAP-SLAM; check association.{]}

{[}31{]} Placed J.A., Strader J., Carrillo H., Atanasov N., Indelman V.,
Carlone L., Castellanos J.A., "A Survey on Active Simultaneous
Localization and Mapping: State of the Art and New Frontiers,"
arXiv:2209.06428, 2022.

{[}32{]} Wu G., Yi T., Fang J., Xie L., Zhang C., Wei W., Liu W., Tian
Q., "4D Gaussian Splatting for Real-Time Dynamic Scene Rendering," Proc.
IEEE/CVF CVPR, 2024. {[}to be verified - cited in text as DynaGS-SLAM;
check association.{]}

{[}33{]} Pan X., Shao W., Zhu Z., Wang Q., Wang T., "GSFF-SLAM: 3D
Semantic Gaussian Splatting for Feature-Field SLAM," arXiv preprint,
2024. {[}to be verified{]}

{[}34{]} Haghighi S., Huang J., Kim B., "OVO-SLAM: Open Vocabulary
Online Simultaneous Localization and Mapping," arXiv preprint, 2024.
{[}to be verified{]}

{[}35{]} Li S., Xue T., Li J., He B., Wu H., "MSGS-SLAM: Multi-Session
Gaussian Splatting SLAM," arXiv preprint, 2024. {[}to be verified{]}

{[}36{]} Li X., Yuan C., Zheng L., Xie J., Zhang J., "NIS-SLAM: Neural
Implicit Semantic RGB-D SLAM for 3D Consistent Scene Understanding,"
arXiv preprint, 2024. {[}to be verified{]}

{[}37{]} Zhang C., Chen Y., Jian M., Hu G., "You Only Group Once:
Efficient Point-Cloud Processing with Token Representation and Relation
Inference Module," Proc. IEEE/RSJ Int. Conf. Intelligent Robots and
Systems (IROS), 2021.

{[}38{]} He M., Guo Y., Ma L., Chen L., Jiang S., "Efficient Point
Transformer for Large-Scale 3D Scene Understanding," arXiv:2405.15827,
2024.

{[}39{]} Guo M.H., Cai J.X., Liu Z.N., Mu T.J., Martin R.R., Hu S.M.,
"PCT: Point Cloud Transformer," Comput. Visual Media, vol. 7, no. 2, pp.
187--199, 2021. DOI: 10.1007/s41095-021-0229-5.

{[}40{]} Ding Z., Yang B., Fan X., Zhou Z., Fan Z., Zhou B., "A Survey
of World Models for Embodied Agents," arXiv:2510.16732, 2025.

{[}41{]} Tordesillas J., Lopez B.T., How J.P., "FASTER: Fast and Safe
Trajectory Planner for Navigation in Unknown Environments," IEEE Trans.
Robot., vol. 38, no. 2, pp. 922--938, 2022. DOI:
10.1109/TRO.2021.3080623. {[}to be verified - cited in context of
risk-aware trajectory planning.{]}

{[}42{]} Ji W., Li J., Bi Q., Liu W., Cheng L., "Point Attention Network
for Semantic Segmentation of 3D Point Clouds," Pattern Recognition, vol.
107, 2020. {[}to be verified - cited in text as SAT3D; check
association.{]}

{[}43{]} Yu X., Tang L., Rao Y., Huang T., Zhou J., Lu J., "Point-BERT:
Pre-Training 3D Point Cloud Transformers with Masked Point Modeling,"
Proc. IEEE/CVF CVPR, 2022. DOI: 10.1109/CVPR52688.2022.01111.

{[}44{]} Kong F., Liu Z., Zhang F., "MARSIM: A Light-weight
Point-realistic Simulator for LiDAR-based UAVs," IEEE Robot. Autom.
Lett., vol. 8, no. 5, 2023. DOI: 10.1109/LRA.2023.3264163.

{[}45{]} Dai Y., Xu H., Ni S., Yin X., Chen B., "NEDS-SLAM: A Novel
Neural Explicit Dense Semantic SLAM Framework Using 3D Gaussian
Splatting," IEEE Robot. Autom. Lett., 2024. {[}to be verified{]}

{[}46{]} Keetha N., Karhade J., Jatavallabhula K.M., Yang G., Scherer
S., Ramanan D., Marques J., "SplaTAM: Splat, Track \& Map 3D Gaussians
for Dense RGB-D SLAM," Proc. IEEE/CVF CVPR, 2024. {[}to be verified -
cited in context of Hi-SLAM {[}5{]}; check association.{]}

{[}47{]} Delmerico J., Scaramuzza D., "A Benchmark Comparison of
Monocular Visual-Inertial Odometry Algorithms for Flying Robots," Proc.
IEEE Int. Conf. Robotics and Automation (ICRA), 2018. DOI:
10.1109/ICRA.2018.8460664. {[}to be verified{]}

{[}48{]} Qin T., Li P., Shen S., "VINS-Mono: A Robust and Versatile
Monocular Visual-Inertial State Estimator," IEEE Trans. Robot., vol. 34,
no. 4, pp. 1004--1020, 2018. DOI: 10.1109/TRO.2018.2853729.

{[}49{]} Alcantarilla P.F., Nuevo J., Bartoli A., "Fast Explicit
Diffusion for Accelerated Features in Nonlinear Scale Spaces," Proc.
British Machine Vision Conference (BMVC), 2013. {[}to be verified -
cited in context of KAZE feature-based SLAM.{]}

{[}50{]} Shan T., Englot B., Meyers D., Wang W., Ratti C., Rus D.,
"LIO-SAM: Tightly-Coupled Lidar Inertial Odometry via Smoothing and
Mapping," Proc. IEEE/RSJ IROS, 2020. DOI:
10.1109/IROS45743.2020.9341176. {[}to be verified{]}

{[}51{]} Zhang J., Singh S., "LOAM: Lidar Odometry and Mapping in
Real-time," Proc. Robotics: Science and Systems (RSS), 2014. {[}to be
verified{]}

{[}52{]} Xu W., Cai Y., He D., Lin J., Zhang F., "Fast-LIO2: Fast Direct
LiDAR-Inertial Odometry," IEEE Trans. Robot., vol. 38, no. 4, pp.
2053--2073, 2022. DOI: 10.1109/TRO.2022.3141876. {[}to be verified{]}

{[}53{]} Nieuwenhuisen M., Droeschel D., Holz D., Behnke S., "Autonomous
Navigation for Micro Aerial Vehicles in Complex GNSS-Denied
Environments," J. Intell. Robot. Syst., vol. 84, no. 1--4, pp. 199--216,
2016. DOI: 10.1007/s10846-015-0274-3.

{[}54{]} Yue W., Nie Y., Song X., Liu G., Wang J., "Semantic-Driven
Visual Navigation for Unmanned Aerial Vehicles in Unknown Environments,"
IEEE Trans. Ind. Electron., vol. 71, no. 11, pp. 13780--13790, 2024.
DOI: 10.1109/tie.2024.3363761.

{[}55{]} Delmerico J., Cieslewski T., Rebecq H., Faessler M., Scaramuzza
D., "Are We Ready for Autonomous Drone Racing? The UZH-FPV Drone Racing
Dataset," Proc. IEEE ICRA, 2019. {[}to be verified - cited in context of
UAV platforms for search-and-rescue missions; check association.{]}

{[}56{]} Brossard M., Barrau A., Bonnabel S., "AI-IMU Dead-Reckoning,"
IEEE Trans. Intell. Veh., vol. 5, no. 4, pp. 585--595, 2020. DOI:
10.1109/TIV.2020.2980758. {[}to be verified - cited in context of
map-aided navigation and dead reckoning.{]}

{[}57{]} Hu J., Ou J., Liu X., Xue B., "A Review of Vision-Based UAV
Localization and Path Planning Based on Optimization for GPS-Denied
Environments," Engineering Research Express, 2025. DOI:
10.1088/2631-8695/ae305b.

{[}58{]} Faessler M., Fontana F., Forster C., Scaramuzza D., "Automatic
Re-Initialization and Failure Recovery for Aggressive Flight with a
Monocular Vision-Based Quadrotor," Proc. IEEE ICRA, 2015. {[}to be
verified - cited in context of UAV GNSS-denied; check association.{]}

{[}59{]} Castano A., "Autonomous UAV Navigation in GPS-Denied
Environments Using LiDAR Point Clouds for Localization and Mapping,"
Proc. AIAA AVIATION Forum, 2023. DOI: 10.2514/6.2023-3626.

\end{document}
